% Section 2 (ontology) of the Milestone 5 report. Owner: Mason Li.
% Written to be \input into the assembled main.tex. Numbers quoted here are
% checked against src/dal/sql/schema.sql and tests/test_semantic_ontology.py.

\section{Proposed Ontology}
\label{sec:ontology}

The accepted Milestone 4 model maps onto one ontology in two steps: each relation
becomes a class, and each column becomes a property. That baseline is applied
without shortcuts, and it accounts for the whole schema rather than a sample of
it. The schema has 51 non-identifier columns: 13 foreign keys and 38 scalar
columns. The keys collapse to 12 distinct forward object properties
(\texttt{hasDimension} serves both the snapshot and the category columns), and
five declared inverses complete the pairs, giving 17 object properties. The
scalar columns collapse to 36 distinct datatype properties, and one further
property, \texttt{recordId}, carries the surrogate row handle the application
still uses, giving 37 datatype properties. Twelve relations therefore yield
twelve classes.

\subsection{Classes, association classes and reused terms}

Two relations are association classes rather than link tables with a payload
attached: \texttt{Holding} carries a fractional interest and its own date, and
\texttt{AllocationEntry} carries the weight that belongs to a snapshot--category
pair. Both stay full classes. Collapsing either into a plain object property
would delete the very quantity the relation exists to hold.

Seven classes specialise FIBO\cite{fibo} or OMG Commons terms, shown in
Table~\ref{tab:alignment}: the portfolio and the holding reuse the FIBO ownership
classes, a snapshot is a \texttt{WeightedBasket} and an entry one of its
weighted constituents, and the scheme and the category reuse the FIBO fund
classification classes. Four published properties are reused as super-properties instead of being
invented: \texttt{weightFraction} is a sub-property of
\texttt{fibo:hasWeight}, both as-of columns are sub-properties of
\texttt{fibo:hasAsOfDate}, which FIBO explicitly distinguishes from a creation
date, and the two containment properties \texttt{hasEntry} and
\texttt{hasHolding} are sub-properties of the OMG Commons
\texttt{hasMember} and \texttt{comprises}. An allocation share is not a
fractional ownership interest, so no FIBO fractional-interest property is
claimed. The alignment is stated by IRI
only; the file carries no \texttt{owl:imports}, so it parses and validates offline
and the reuse stays an auditable claim rather than a network dependency.

% FIBO alignment table for Section 2. Owner: Mason Li.
% Local terms are in the etf namespace; inherited module paths are in
% milestones/05_semantic/docs/fibo_alignment.md.
\begin{table}[t]
\centering
\scriptsize
\caption{The reuse boundary. Seven classes specialise FIBO terms; five classes and
the observation attributes use project-specific definitions.}
\label{tab:alignment}
\begin{tabular}{@{}l@{}}
\toprule
\multicolumn{1}{@{}l}{\itshape reused (subsumption claim)}\\
etf:ETF $\sqsubseteq$ fibo:CollectiveInvestmentVehicle \\
etf:Portfolio $\sqsubseteq$ fibo:Portfolio \\
etf:Holding $\sqsubseteq$ fibo:Holding \\
etf:AllocationSnapshot $\sqsubseteq$ fibo:WeightedBasket \\
etf:AllocationEntry $\sqsubseteq$ fibo:WeightedBasketConstituent \\
etf:ClassificationScheme $\sqsubseteq$ fibo:FundClassificationScheme \\
etf:AllocationCategory $\sqsubseteq$ fibo:FundClassification \\
\midrule
\multicolumn{1}{@{}l}{\itshape extensions (no FIBO parent)}\\
Customer, ETFProvider (publisher role declined), ETFIdentifier, \\
AllocationDimension, SourceDocument (capture provenance) \\
\bottomrule
\end{tabular}
\end{table}


\subsection{Identity and context}

Identity is minted from the columns that already separate two records, never from
the surrogate \texttt{Id}, which is an artefact of row storage. Customers, providers,
dimensions and identifiers get \texttt{owl:hasKey}. Composite keys that name a
parent---a portfolio under a customer, a snapshot under a fund---do not: a
\texttt{hasKey} over an object property asserts individual equality across
parents, which is stronger than per-parent uniqueness. Those keys become URI
templates that mirror the constraint, one path segment per key component in key
order, with an explicit duplicate check beside them.

The category is the one class whose identity is not a single column. The schema
keys it by (dimension, context key, label). The context key is a prefix plus a
JSON list: either the classification authority, name and version, or --- when the
edition is unpublished --- the source URL, retrieval time, locator and dimension.
The graph reproduces that construction, and a parity test asserts that the graph
helper and the relational helper return byte-identical strings, so the two layers
cannot drift into disagreeing about which labels are the same label.

Provenance is structural, not annotational. Every snapshot keeps its source
document, its as-of date and the locator inside the capture that the weights were
read from, and the capture itself is retained by reference. An unresolved
classification edition keeps its version unrecorded instead of being assigned a
fabricated one, and a second unversioned edition for the same authority is
refused. Under the synthetic demonstration the retained file is generated, so the
file-existence check proves only that a capture exists---never that it supports
the numbers.

\subsection{What OWL states, and what it cannot}

The ontology states the schema's shape: domains and ranges, exactly-one and
at-least-one restrictions on participation, at-most-one scheme per category,
functional properties, four keys, and disjointness between snapshot, entry and
category. One axiom is worth singling out. The four schema triggers keep an
entry's category dimension equal to its snapshot's dimension; the graph restates
that as the property chain
\texttt{hasEntry}~$\circ$~\texttt{usesCategory}~$\circ$~\texttt{hasDimension}
$\sqsubseteq$~\texttt{hasDimension}.

But an open-world ontology\cite{owl2profiles} does not forbid what it does not
know, and it cannot add. A missing basis note is unknown, not invalid; a total of 0.94 against
\texttt{Complete} is not an axiom violation; a closed value set such as
\{Sector, Geography\} needs enumerated individuals, which would drag instance data
into the schema file. The dimension chain entails agreement but still cannot
\emph{flag} a disagreement, because without an \texttt{AllDifferent} assertion on
the dimension individuals the two dimensions may simply be the same individual.
Table~\ref{tab:split} records the split we ship: axioms for what must hold,
explicit checks for aggregation, closed sets, filesystem facts and composite
uniqueness. The validator module, built on RDFLib~\cite{rdflib}, publishes
both lists and a CLI, so the split can be reviewed without reading the code.

% OWL-versus-explicit rule split for Section 2. Owner: Mason Li.
\begin{table}[t]
\centering
\footnotesize
\caption{Where each accepted constraint is enforced. Four groups cannot be
axioms: aggregates, closed value sets, filesystem facts, and uniqueness that names
a parent.}
\label{tab:split}
\begin{tabular}{@{}>{\raggedright\arraybackslash}p{0.46\columnwidth}>{\raggedright\arraybackslash}p{0.50\columnwidth}@{}}
\toprule
Expressed as OWL/RDFS axioms & Requires an explicit check \\
\midrule
domains and ranges of all 54 properties & weights in $[0,1]$ \\
exactly-one participation, twelve functional properties & \texttt{Complete} total $=1\pm0.005$ (a sum) \\
at-most-one scheme per category & entry dimension $=$ snapshot dimension \\
\texttt{owl:hasKey} on four classes & closed sets: Sector/Geography, ISIN/Ticker, Complete/Partial/Unknown \\
property chain for dimension agreement & ISIN and ticker lexical shapes \\
disjointness of snapshot/entry/category & NOT NULL slots (open world reads absence as unknown) \\
FIBO subsumptions & context key against its scheme edition \\
 & composite keys that name a parent \\
 & capture file present inside the capture root \\
\bottomrule
\end{tabular}
\end{table}


The split is also the argument for the extension. FIBO's
\texttt{FundPortfolioInvestmentPolicy} describes how a portfolio is \emph{intended}
to be allocated; a snapshot records what a provider \emph{disclosed} on a date,
sometimes with partial coverage. Subsuming one under the other would license
inferences about managerial intent that the evidence does not support, so the
reuse stops at the basket and the constituent and the observation is modelled as
an extension with coverage status, basis note and locator. We likewise decline to
type the provider as FIBO's \texttt{Publisher} role: one relation plays both the
issuer and the publisher role, and the correct FIBO-shaped treatment would be two
role individuals over one party, a modelling change outside this milestone.
